\QHeading{第1题　用简便方法计算下面各题。}
\par\noindent


\CalcBlock[方法一：简便计算]{1}{\begin{aligned}[t]&\phantom{={}}\color{black}676+187+124\\&\color{answerred}=676+124+187\\&\color{answerred}=(676+124)+187\\&\color{answerred}=800+187\\&\color{answerred}=987\end{aligned}}
\par
\MethodBlock{方法二：按运算顺序直接计算（对照）}{本法只作结果核对；本题要求简便计算，前面的简便解法仍保留。}{\begin{aligned}[t]&\phantom{={}}676+187+124\\
&=863+124\\
&=987\end{aligned}}


\CalcBlock[方法一：简便计算]{2}{\begin{aligned}[t]&\phantom{={}}\color{black}236+183+117\\&\color{answerred}=236+(183+117)\\&\color{answerred}=236+300\\&\color{answerred}=536\end{aligned}}
\MethodBlock{方法二：按运算顺序直接计算（对照）}{本法只作结果核对；本题要求简便计算，前面的简便解法仍保留。}{\begin{aligned}[t]&\phantom{={}}236+183+117\\
&=419+117\\
&=536\end{aligned}}

\CalcBlock[方法一：简便计算]{3}{\begin{aligned}[t]&\phantom{={}}\color{black}557-(257+158)\\&\color{answerred}=557-257-158\\&\color{answerred}=300-158\\&\color{answerred}=142\end{aligned}}
\MethodBlock{方法二：按运算顺序直接计算（对照）}{本法只作结果核对；本题要求简便计算，前面的简便解法仍保留。}{\begin{aligned}[t]&\phantom{={}}557-(257+158)\\
&=557-415\\
&=142\end{aligned}}

\CalcBlock[方法一：简便计算]{4}{\begin{aligned}[t]&\phantom{={}}\color{black}474-198-174\\&\color{answerred}=474-174-198\\&\color{answerred}=(474-174)-198\\&\color{answerred}=300-198\\&\color{answerred}=102\end{aligned}}
\MethodBlock{方法二：先合并要减去的部分}{连续减去几个数，可以先把这些数相加再减。}{\begin{aligned}[t]&\phantom{={}}474-198-174\\
&=474-(198+174)\\
&=474-372\\
&=102\end{aligned}}

\MethodBlock{方法三：按运算顺序直接计算（对照）}{本法只作结果核对；本题要求简便计算，前面的简便解法仍保留。}{\begin{aligned}[t]&\phantom{={}}474-198-174\\
&=276-174\\
&=102\end{aligned}}

\CalcBlock[方法一：简便计算]{5}{\begin{aligned}[t]&\phantom{={}}\color{black}860-125-175\\&\color{answerred}=860-(125+175)\\&\color{answerred}=860-300\\&\color{answerred}=560\end{aligned}}
\MethodBlock{方法二：按运算顺序直接计算（对照）}{本法只作结果核对；本题要求简便计算，前面的简便解法仍保留。}{\begin{aligned}[t]&\phantom{={}}860-125-175\\
&=735-175\\
&=560\end{aligned}}

\CalcBlock[方法一：简便计算]{6}{\begin{aligned}[t]&\phantom{={}}\color{black}548-298\\&\color{answerred}=548-(300-2)\\&\color{answerred}=548-300+2\\&\color{answerred}=248+2\\&\color{answerred}=250\end{aligned}}
\MethodBlock{方法二：补差法}{从减数先补到一个整数，再补到被减数，两段差相加。}{\begin{aligned}[t]&\phantom{={}}548-298\\
&=(300-298)+(548-300)\\
&=2+248\\
&=250\end{aligned}}

\CalcBlock[方法一：简便计算]{7}{\begin{aligned}[t]&\phantom{={}}\color{black}99\times23\\&\color{answerred}=(100-1)\times23\\&\color{answerred}=100\times23-1\times23\\&\color{answerred}=2300-23\\&\color{answerred}=2277\end{aligned}}
\MethodBlock{方法二：按数位拆分计算}{把一个因数拆成整十、整百数与余下部分之和。}{\begin{aligned}[t]&\phantom{={}}99\times 23\\
&=99\times (20+3)\\
&=99\times 20+99\times 3\\
&=1980+297\\
&=2277\end{aligned}}

\CalcBlock[方法一：简便计算]{8}{\begin{aligned}[t]&\phantom{={}}\color{black}101\times36\\&\color{answerred}=(100+1)\times36\\&\color{answerred}=100\times36+1\times36\\&\color{answerred}=3600+36\\&\color{answerred}=3636\end{aligned}}
\MethodBlock{方法二：按数位拆分计算}{把一个因数拆成整十、整百数与余下部分之和。}{\begin{aligned}[t]&\phantom{={}}101\times 36\\
&=101\times (30+6)\\
&=101\times 30+101\times 6\\
&=3030+606\\
&=3636\end{aligned}}

\CalcBlock[方法一：简便计算]{9}{\begin{aligned}[t]&\phantom{={}}\color{black}254-98\\&\color{answerred}=254-(100-2)\\&\color{answerred}=254-100+2\\&\color{answerred}=154+2\\&\color{answerred}=156\end{aligned}}
\MethodBlock{方法二：补差法}{从减数先补到一个整数，再补到被减数，两段差相加。}{\begin{aligned}[t]&\phantom{={}}254-98\\
&=(100-98)+(254-100)\\
&=2+154\\
&=156\end{aligned}}

\par\noindent

\QHeading{第2题　解决问题。}

\Problem{① 文文带着30元钱去买文具，买钢笔用去9.8元，买文具盒比买钢笔多用去4.2元。文文还剩下多少钱？}
\par


\MethodTitle{方法一}
\Answer{{\Formula{\begin{aligned}[t]&\phantom{={}}30-9.8-(9.8+4.2)\\&=30-9.8-14\\&=20.2-14\\&=6.2\text{（元）}\end{aligned}}}\par 答：文文还剩6.2元。}

\MethodTitle{方法二：先按两件都花9.8元，再补文具盒的差额}
\Answer{两件商品一共是两个9.8元再加4.2元，剩余用总钱数减去。}
\Answer{\Formula{\begin{aligned}[t]&\phantom{={}}30-(9.8\times 2+4.2)\\
&=30-(19.6+4.2)\\
&=30-23.8\\
&=6.2\text{（元）}\end{aligned}}}
\Answer{答：文文还剩6.2元。}
\Problem{② 一根绳子长10米，第一次用去3.5米，第二次又用去了一些，这时剩下的部分比第一次用去的短0.6米。第二次用去多少米绳子？}

\MethodTitle{方法一：简便计算}
\Answer{{\Formula{\begin{aligned}[t]&\phantom{={}}10-3.5-(3.5-0.6)\\&=10-3.5-3.5+0.6\\&=10-(3.5+3.5)+0.6\\&=10-7+0.6\\&=3+0.6\\&=3.6\text{（米）}\end{aligned}}}\par 答：第二次用去3.6米。}

\MethodTitle{方法二：先把第一次用量与最后剩余量合并}
\Answer{剩余比第一次用量少0.6米，两部分合计是两个3.5米再减0.6米。}
\Answer{\Formula{\begin{aligned}[t]&\phantom{={}}10-(3.5\times 2-0.6)\\
&=10-(7-0.6)\\
&=10-6.4\\
&=3.6\text{（米）}\end{aligned}}}
\Answer{答：第二次用去3.6米绳子。}
